%% Included in decks: slidemage-manual
%%

% SlideMage manual, chapter "Assembling decks": the \slide family.
\begin{standardframe}{Six ways to put a slide into a deck}
	{\small A deck lists its slides with one command each. Mostly that is
	\SMCmd{slide}; the others give a deck an old version of a slide, a slide fully
	revealed, or a slide in some deck versions only.\par}
	\smallskip
	\begin{columns}[T,onlytextwidth]
		\begin{column}{0.6\linewidth}
			\footnotesize
			\renewcommand{\arraystretch}{0.9}
			\begin{gothamtable}[gap=8pt]{ll}
				\toprule
				Call & What the deck gets \\
				\midrule
				\SMCmd{slide} & the live \SMCode{current.tex} \\
				\SMCmd{slideversion} & a frozen version, pinned \\
				\SMCmd{slidefinal} & fully revealed, in every version \\
				\SMCmd{conditionalslide} & only in the listed versions \\
				\SMCmd{slidestep} & reveal steps chosen per version \\
				\SMCmd{omitslide} & the older way (\SMCmd{studentversion}) \\
				\bottomrule
			\end{gothamtable}
		\end{column}
		\begin{column}{0.37\linewidth}
			\SMCodeFile[basicstyle=\SMCodeFont\fontsize{6.2}{7.8}\selectfont]{snippets/1u4-1.tex}
			\smallskip
			{\small The deck editor writes them for you.\par}
		\end{column}
	\end{columns}
	\smallskip
	{\small
	\begin{alertblock}{Slides with a course date}
		Each deck pins its own edition with \SMCmd{slideversion}; a plain \SMCmd{slide} would let one course's edit change all.
	\end{alertblock}}
\end{standardframe}
